% id: 0484
% book: RPM 기하 · 05 벡터의 연산
% section: 유형 05 평면도형에서 벡터의 연산
% figure: crop:fig-0484.png
% answer: $4\vec{a}+3\vec{b}$
% answer_source: 답지
% variant_level: 0 (원본 전사)
% 이 파일은 build.mjs 가 items.json 에서 만든다. 고칠 때는 items.json 을 고친다.
\smash{\raisebox{1.5mm}{\dmkichul{RPM 기하 0484번 · 77쪽 · 서술형}}}
\noindent\begin{minipage}[t]{0.58\linewidth}\vspace{0pt}\raggedright 오른쪽 그림과 같이 합동인 정육각형 $2$개가 한 변을 공유하고 있다. $\overrightarrow{\pt{OA}}=\vec{a}$, $\overrightarrow{\pt{OB}}=\vec{b}$라 할 때, $\overrightarrow{\pt{BQ}}+\overrightarrow{\pt{PR}}$를 $\vec{a}$, $\vec{b}$로 나타내시오.\end{minipage}\hfill\begin{minipage}[t]{0.4\linewidth}\vspace{0pt}\centering \setlength{\dmfigw}{113.2pt}\ifdim\dmfigw>\linewidth\setlength{\dmfigw}{\linewidth}\fi\includegraphics[width=\dmfigw]{figures/fig-0484.png}\end{minipage}\par
